\scanpage{81}
\mspage{80}

\uline{Corollaire} Soient $F,G$ avec $F\ll G$. Pour tout $X\in\bar F$, soit $\varphi(X)\in\bar G$ la plus petite \deleted{multi}strate \inserted{$Y$} de $G$ telle que $\deleted{\varphi(X)\ll}\,\unsure{.}\,X\ll Y$. L'application
\[
 \varphi:\bar F\longrightarrow\bar G
\]
ainsi obtenue est \uline{croissante} (trivial).

Jusqu'à \unsure{maintnt|maintenant}, les axiomes $C5$, $C6$, $C7$ introduits pour $\preceq$ \deleted{\unsure{pour}} sont tautologiques, si on prend pour $\preceq$ l'égalité, \inserted{alors $\ll$ est $\leq$, \unsure{i.e.|scil.} l'inclusion de figures.} Il serait temps d'en venir à des propriétés \deleted{telles que} plus spécifiques !

\uline{14 juin}

\margin{Je vais poser $\widetilde F$ ald $\bar F$ pour l'ens des $X\in M$ tels que $X\leq F$ (pour éviter confusion avec notion d'\unsure{adhérence|adhésion})}

\deleted{\unsure{\uuline{Proposition}|\uline{Proposition}}} \inserted{Corollaire de prop. \unsure{pag.|page} 76} Cf l'\uline{ancien} $C4$ \deleted{\unsure{pa|pag|p}} page 47)

Soient $G'\deleted{\unsure{\leq}}\ll G\deleted{\unsure{\leq}}\preceq F$ alors $\widetilde{G'}\cap\widetilde F\subset\widetilde G$
\inserted{(donc $\widetilde G'\cap\widetilde F=\widetilde{G'}\cap\widetilde G$)}

En d'autres termes, toute \inserted{(multi)}strate \deleted{\unsure{de}} \deleted{\unsure{la} \unsure{figure}} \inserted{$X$ \unsure{(}qui raffine la sous-figure $G$ de $F$ \unsure{et|)}} est \inserted{(multi)}strate de $F$ est \inserted{(multi)}strate de $G$.

\deleted{\deleted{Soit $X$ strate de $G'$ qui est strate de $F$}}

\deleted{\deleted{\unsure{D}} \unsure{Dém.|Dem.} $X$ est une strate d'un raffinement $\bar G$ de $G$, \inserted{lequel est} induit par un raffinement $\deleted{\unsure{G}}\unsure{\bar F|F}$ de $F$ ($C6$)}
\diagram{\begin{tikzcd}[ampersand replacement=\&,row sep=0.05em,column sep=0.16em]
\deleted{G} \& \deleted{\unsure{=|\leq}} \& \deleted{F} \& \deleted{\geq} \& \deleted{X}\\
\deleted{\rotatebox{90}{$\preceq$}} \& \& \deleted{\rotatebox{90}{$\preceq$}} \& \& \deleted{\unsure{;}}\\
\deleted{\bar G} \& \deleted{\leq} \& \deleted{\unsure{.|\cdot}\bar F} \& \& \\
\deleted{\rotatebox{90}{$\leq$}} \& \& \& \& \\
\deleted{X} \& \& \& \&
\end{tikzcd}}
\margin{\unsure{hi hi hi|si si si|hé hé hé|q e.d.|q f.d.} $\longrightarrow$}
\deleted{Supposons $X\leq F$, \deleted{et considérons les \unsure{images|image} \unsure{inverses|inverse}} \inserted{l'image inverse de}}
